% ======================================================================
%  Định dạng báo cáo — theo mẫu Tiểu luận 03 (Times 13pt, giãn dòng 1.5)
% ======================================================================
\usepackage[a4paper,left=3cm,right=2cm,top=2.5cm,bottom=2.5cm,headheight=14pt]{geometry}
\usepackage{fontspec}
\usepackage{unicode-math}
\setmainfont{TeX Gyre Termes}[Scale=1.0833]          % 12pt × 1,0833 ≈ 13pt
\setmathfont{TeX Gyre Termes Math}[Scale=1.0833]
\setmonofont{SourceCodePro}[Extension=.otf,UprightFont=*-Regular,BoldFont=*-Bold,ItalicFont=*-RegularIt,BoldItalicFont=*-BoldIt,Scale=0.82]
\usepackage{polyglossia}
\setdefaultlanguage{vietnamese}
\setotherlanguage{english}  % để biblatex dùng chuỗi tiếng Anh (đã Việt hóa) trong danh mục tài liệu

\usepackage{setspace}
\setstretch{1.4}
\setlength{\parindent}{1cm}
\setlength{\parskip}{3pt}
\usepackage{indentfirst}
\usepackage{microtype}
\usepackage{ragged2e}
\setlength{\emergencystretch}{3em}

\usepackage{amsmath,amsthm}
\usepackage{graphicx}
\graphicspath{{figures/}}
\usepackage{float}
\usepackage{booktabs,longtable,tabularx,multirow,array,makecell}
\renewcommand{\arraystretch}{1.18}
\usepackage{xcolor}
\definecolor{codebg}{HTML}{F4F5F7}
\definecolor{accent}{HTML}{1D4ED8}
\usepackage[hidelinks]{hyperref}
\usepackage{xurl}
\usepackage[htt]{hyphenat}  % cho phép ngắt dòng trong \texttt (tên hàm, đường dẫn dài)
\usepackage{adjustbox}
\usepackage{enumitem}
\setlist{itemsep=1pt,topsep=3pt,parsep=0pt,leftmargin=1.2cm}
\usepackage{tikz}
\usetikzlibrary{positioning,arrows.meta,shapes.geometric,calc,fit,backgrounds}

% ---------------------------------------------------------------- tiêu đề chương, mục
\usepackage{titlesec}
\titleformat{\chapter}[block]{\centering\bfseries\fontsize{15}{19}\selectfont}
  {\MakeUppercase{\chaptertitlename}\ \thechapter.}{0.5em}{\MakeUppercase}
\titlespacing{\chapter}{0pt}{-10pt}{18pt}
\titleformat{\section}{\bfseries\fontsize{13.5}{17}\selectfont}{\thesection.}{0.5em}{}
\titleformat{\subsection}{\bfseries\itshape\fontsize{13}{16}\selectfont}{\thesubsection.}{0.5em}{}
\titleformat{\subsubsection}{\itshape\fontsize{13}{16}\selectfont}{\thesubsubsection.}{0.5em}{}
\titlespacing{\section}{0pt}{10pt}{4pt}
\titlespacing{\subsection}{0pt}{8pt}{3pt}
\titlespacing{\subsubsection}{0pt}{6pt}{2pt}
\setcounter{secnumdepth}{3}
\setcounter{tocdepth}{2}
\renewcommand{\chaptertitlename}{Chương}

% chương không đánh số (Mở đầu, Kết luận, TLTK) vẫn vào mục lục
\newcommand{\unnumberedchapter}[1]{%
  \chapter*{#1}\addcontentsline{toc}{chapter}{#1}\markboth{#1}{#1}}

% ---------------------------------------------------------------- mục lục
\usepackage{tocloft}
\renewcommand{\contentsname}{\centering MỤC LỤC}
\renewcommand{\listfigurename}{\centering DANH MỤC HÌNH VẼ}
\renewcommand{\listtablename}{\centering DANH MỤC BẢNG BIỂU}
\renewcommand{\cftchapfont}{\bfseries}
\renewcommand{\cftchappagefont}{\bfseries}
\renewcommand{\cftchappresnum}{CHƯƠNG }
\setlength{\cftchapnumwidth}{6.6em}
\renewcommand{\cftchapaftersnum}{.}
\renewcommand{\cftsecaftersnum}{.}
\renewcommand{\cftsubsecaftersnum}{.}
\renewcommand{\cftchapleader}{\cftdotfill{\cftdotsep}}
\renewcommand{\cftfigpresnum}{Hình }
\renewcommand{\cfttabpresnum}{Bảng }
\setlength{\cftfignumwidth}{5.2em}
\setlength{\cfttabnumwidth}{5.2em}
\renewcommand{\cftfigaftersnum}{.}
\renewcommand{\cfttabaftersnum}{.}

% ---------------------------------------------------------------- chú thích hình, bảng
\usepackage{caption}
\DeclareCaptionLabelSeparator{dotspace}{.\ }
\captionsetup{labelsep=dotspace,font={small},justification=centering,singlelinecheck=false}
\captionsetup[figure]{name=Hình,font={small,it},skip=4pt}
\captionsetup[table]{name=Bảng,font={small,bf},position=top,skip=4pt}
\renewcommand{\thefigure}{\thechapter.\arabic{figure}}
\renewcommand{\thetable}{\thechapter.\arabic{table}}

% ---------------------------------------------------------------- đoạn mã
\usepackage[newfloat]{minted}
\setminted{fontsize=\footnotesize,breaklines,bgcolor=codebg,frame=none,tabsize=4,
           baselinestretch=1.05,style=friendly,autogobble}
\SetupFloatingEnvironment{listing}{name=Đoạn mã,placement=H}
\captionsetup[listing]{font={small,bf},position=top,justification=raggedright,skip=2pt}
\renewcommand{\thelisting}{\thechapter.\arabic{listing}}
\newenvironment{code}[1]{\begin{listing}[H]\caption{#1}}{\end{listing}}

% ---------------------------------------------------------------- trang
\usepackage{fancyhdr}
\pagestyle{fancy}
\fancyhf{}
\fancyfoot[C]{\thepage}
\renewcommand{\headrulewidth}{0pt}
\fancypagestyle{plain}{\fancyhf{}\fancyfoot[C]{\thepage}\renewcommand{\headrulewidth}{0pt}}

% ---------------------------------------------------------------- tài liệu tham khảo
\usepackage[backend=biber,style=ieee,sorting=none,maxnames=4,minnames=1,language=english,autolang=none]{biblatex}
\addbibresource{refs.bib}
\DefineBibliographyStrings{english}{andothers={và cộng sự},and={và},in={trong},edition={ấn bản},volume={tập},jourvol={tập},number={số},pages={tr.},page={tr.}}

% ---------------------------------------------------------------- tiện ích
\newcommand{\eq}[1]{(\ref{#1})}
\newcommand{\figref}[1]{Hình~\ref{#1}}
\newcommand{\tabref}[1]{Bảng~\ref{#1}}
\newcommand{\coderef}[1]{Đoạn mã~\ref{#1}}
\newcommand{\weburl}[1]{{\small\url{#1}}}
\newcommand{\api}[1]{\texttt{\small #1}}
\newcommand{\figw}[3][\linewidth]{%
  \begin{figure}[H]\centering\includegraphics[width=#1]{#2}\caption{#3}\label{fig:#2}\end{figure}}
\newcolumntype{L}[1]{>{\raggedright\arraybackslash}p{#1}}
\newcolumntype{C}[1]{>{\centering\arraybackslash}p{#1}}
\newcolumntype{Y}{>{\raggedright\arraybackslash}X}
\newtheorem*{dinhnghia}{Định nghĩa}

% tên miền triển khai (Cloudflare Tunnel) — đổi ở đây là đổi toàn bộ báo cáo
\newcommand{\deploydomain}{agent.nits.io.vn}
